\section{A compact Euler integral and its unavoidable logarithmic
contamination}
\label{sec:compact-integral-contamination}

The arithmetic framework of this section---Beukers--Sondow double
integrals and Nesterenko-type rational functions producing linear forms in
\(1,\gamma\), and logarithms---goes back at least to Sondow
\cite{Sondow2003Criteria,Sondow2009Hypergeometric}.  We make no priority
claim for that framework.  We instead isolate the exact asymmetric
contiguous family needed here, correct its asymptotic and height
normalization, and prove that its logarithmic contamination is strictly
positive for every index.

\subsection{The floor-weighted rational transform}

For a rational function \(g\) which is integrable against the indicated
weight, define
\begin{equation}\label{eq:floor-transform}
 \mathcal T(g)=\int_0^\infty(\lfloor u\rfloor+1)g(u)\,du.
\end{equation}
Write \(H_m^{(r)}=\sum_{j=1}^m j^{-r}\), with \(H_0^{(r)}=0\), and
\(H_m=H_m^{(1)}\).

\begin{proposition}[Complete pole-order transducer]
\label{prop:complete-pole-transducer}
Suppose that \(g\in\Qbar(u)\) is \(O(u^{-3})\) at infinity, all its poles
belong to \(\{-1,-2,\ldots\}\), and
\[
 g(u)=\sum_{p\ge1}\sum_{r=1}^{m_p}
       \frac{c_{p,r}}{(u+p)^r}
\]
is its finite partial-fraction expansion.  Then
\begin{align}
 \mathcal T(g)
 ={}&\sum_p c_{p,1}
 \left(\log((p-1)!)+p+1-\frac12\log(2\pi)\right)
 \notag\\
 &+\sum_p c_{p,2}\left(\gamma-1-H_{p-1}\right)
 \label{eq:complete-pole-transducer}\\
 &+\sum_p\sum_{r=3}^{m_p}
 \frac{c_{p,r}}{r-1}
 \left(\zeta(r-1)-H_{p-1}^{(r-1)}\right).
 \notag
\end{align}
In particular, simple poles carry factorial logarithms, double poles carry
\(\gamma\), and poles of order \(r\ge3\) carry \(\zeta(r-1)\).
\end{proposition}

\begin{proof}
For a simple pole, summation over the intervals \([j,j+1]\) gives
\[
 \sum_{j=0}^J(j+1)
 \log\frac{j+p+1}{j+p}
 =(J+1)\log(J+p+1)-\log\Gamma(J+p+1)+\log\Gamma(p).
\]
Stirling's formula isolates the finite part
\(\log\Gamma(p)+p+1-\frac12\log(2\pi)\).
For a double pole the corresponding finite part is
\[
 \frac1p-1-\psi(p+1)=\gamma-1-H_{p-1}.
\]
For \(r\ge3\), absolute convergence and summation by parts give
\begin{align*}
 &\sum_{j\ge0}(j+1)\int_j^{j+1}\frac{du}{(u+p)^r}\\
 &\hspace{2cm}=\frac1{r-1}\sum_{j\ge0}(j+p)^{1-r}
 =\frac{\zeta(r-1,p)}{r-1}.
\end{align*}
Since \(p\) is a positive integer,
\(\zeta(r-1,p)=\zeta(r-1)-H_{p-1}^{(r-1)}\).

The hypothesis \(g(u)=O(u^{-3})\) gives
\[
 \sum_pc_{p,1}=0,
 \qquad
 \sum_pc_{p,2}=\sum_ppc_{p,1}.
\]
These identities cancel the linear and logarithmic divergences of the
simple- and double-pole partial sums, so the sum of their finite parts is
exactly \eqref{eq:complete-pole-transducer}.
\end{proof}

\begin{remark}[Necessary correction of scope]
A product of Beta factors can have poles of multiplicity greater than two.
In that situation a partial-fraction expansion containing only simple and
double poles is incomplete; the last line of
\eqref{eq:complete-pole-transducer} is indispensable.  Restricting to
products with pole multiplicity at most two recovers the
\((1,\gamma,\log)\)-valued special case.
\end{remark}

There is also a one-variable endpoint obstruction which explains why the
compact construction below is genuinely two-dimensional.

\begin{lemma}[Endpoint removal loses \(\gamma\)]
\label{lem:one-variable-endpoint-no-go}
Let \(P(x)=\sum_{k=0}^d a_kx^k\in\Qbar[x]\) satisfy \(P(1)=0\).  Then
\begin{equation}\label{eq:endpoint-no-go-exact}
 \int_0^1P(x)\left(\frac1{1-x}+\frac1{\log x}\right)dx
 =\int_0^1\frac{P(x)}{1-x}\,dx
  +\sum_{k=0}^da_k\log(k+1).
\end{equation}
The first term on the right is algebraic; in particular no \(\gamma\)-term
survives.
\end{lemma}

\begin{proof}
The first integrand is a polynomial because \(P(1)=0\).  The same condition
allows one to write \(P(x)=\sum_ka_k(x^k-1)\), and Frullani's identity
\[
 \int_0^1\frac{x^k-1}{\log x}\,dx=\log(k+1)
\]
proves \eqref{eq:endpoint-no-go-exact}.
\end{proof}

\subsection{The asymmetric compact family}

For \(n\ge1\), set
\begin{equation}\label{eq:asymmetric-Jn}
 J_n=\int_0^1\!\int_0^1
 \frac{[x(1-x)y(1-y)]^n(1-x)}
 {(1-xy)(-\log(xy))}\,dx\,dy.
\end{equation}
Using
\(1/(-\log(xy))=\int_0^\infty(xy)^t\,dt\) and
\((1-xy)^{-1}=\sum_{j\ge0}(xy)^j\), Tonelli's theorem gives
\begin{equation}\label{eq:Jn-floor-beta}
 J_n=\mathcal T(g_n),\qquad
 g_n(u)=B(u+n+1,n+2)B(u+n+1,n+1).
\end{equation}
Put \(L=n+1\), \(M=2n+2\), and \(C=n!(n+1)!\).  Then
\begin{equation}\label{eq:gn-rational-form}
 g_n(u)=\frac{C}{(u+M)\prod_{p=L}^{M-1}(u+p)^2}.
\end{equation}

\begin{theorem}[Exact coefficients and strict logarithmic contamination]
\label{thm:Jn-exact-strict-contamination}
The partial-fraction expansion of \(g_n\) is
\[
 g_n(u)=\sum_{p=L}^{M}\frac{\alpha_p}{u+p}
       +\sum_{p=L}^{M-1}\frac{\beta_p}{(u+p)^2},
\]
where, for \(0\le k\le n\),
\begin{align}
 \beta_{L+k}
 &=\binom nk\binom{n+1}{k},
 \label{eq:Jn-beta}\\
 \alpha_{L+k}
 &=\beta_{L+k}\left(
 2H_k-2H_{n-k}-\frac1{n+1-k}\right),
 \label{eq:Jn-alpha}\\
 \alpha_M&=\frac1{n+1}.
 \label{eq:Jn-alpha-last}
\end{align}
Consequently
\begin{equation}\label{eq:Jn-exact-linear-form}
 J_n=\binom{2n+1}{n}\gamma+r_n+\mathcal L_n,
\end{equation}
where
\begin{align}
 r_n
 &=\sum_{p=L}^{M}\alpha_p(p+1)
   +\sum_{p=L}^{M-1}\beta_p
      \left(\frac1p-1-H_p\right)\in\Q,
 \label{eq:Jn-rational-part}\\
 \mathcal L_n
 &=\sum_{p=L}^{M}\alpha_p\log((p-1)!)
   =\sum_{j=L}^{M-1}T_j\log j,
 \qquad
 T_j=\sum_{p>j}\alpha_p.
 \label{eq:Jn-log-part}
\end{align}
Every proper tail is strictly positive:
\begin{equation}\label{eq:Jn-positive-tails}
 T_j>0\qquad(L\le j\le M-1).
\end{equation}
In particular
\begin{equation}\label{eq:Jn-log-positive}
 \boxed{\mathcal L_n>0\quad\text{for every }n\ge1.}
\end{equation}
Thus this entire compact family is never logarithm-free.
\end{theorem}

\begin{proof}
At the double pole \(p=L+k\), evaluation of
\((u+p)^2g_n(u)\) at \(u=-p\) gives
\[
 \beta_{L+k}
 =\frac{n!(n+1)!}
 {(n+1-k)(k!)^2((n-k)!)^2}
 =\binom nk\binom{n+1}{k}.
\]
Logarithmic differentiation of the same expression gives
\[
 \frac{\alpha_{L+k}}{\beta_{L+k}}
 =-\frac1{n+1-k}
  -2\sum_{\substack{0\le j\le n\\j\ne k}}\frac1{j-k}
 =2H_k-2H_{n-k}-\frac1{n+1-k},
\]
which proves \eqref{eq:Jn-alpha}.  The simple pole at \(M\) gives
\eqref{eq:Jn-alpha-last}.  Since \(g_n(u)=O(u^{-2n-3})\), comparison at
infinity gives
\begin{equation}\label{eq:Jn-alpha-sum-zero}
 \sum_{p=L}^{M}\alpha_p=0.
\end{equation}

Proposition~\ref{prop:complete-pole-transducer} now yields
\eqref{eq:Jn-exact-linear-form}--\eqref{eq:Jn-log-part}; the
\(\log(2\pi)\)-term disappears by
\eqref{eq:Jn-alpha-sum-zero}.  Vandermonde's identity gives the
\(\gamma\)-coefficient:
\[
 \sum_{k=0}^n\beta_{L+k}
 =\sum_{k=0}^n\binom nk\binom{n+1}{k}
 =\binom{2n+1}{n}.
\]

It remains to prove the strict obstruction.  Write
\[
 b_k=2H_k-2H_{n-k}-\frac1{n+1-k}.
\]
Then
\[
 b_{k+1}-b_k
 =\frac2{k+1}+\frac1{n-k}+\frac1{n+1-k}>0,
\]
while \(b_0<0\) and \(b_n=2H_n-1>0\).  Since every
\(\beta_{L+k}\) is positive, the sequence
\(\alpha_L,\ldots,\alpha_{M-1}\) changes sign at most once, from negative
to positive, and \(\alpha_M>0\).  Together with the zero total sum
\eqref{eq:Jn-alpha-sum-zero}, this implies that every proper tail \(T_j\)
is positive.  Formula \eqref{eq:Jn-log-part} then gives
\eqref{eq:Jn-log-positive}.
\end{proof}

\subsection{Corrected asymptotic and integer height}

\begin{theorem}[Corrected asymptotic and integerization]
\label{thm:Jn-asymptotic-integerization}
As \(n\to\infty\),
\begin{equation}\label{eq:Jn-correct-asymptotic}
 \boxed{\displaystyle
 J_n\sim\frac{\pi}{12\log2}\,\frac{16^{-n}}n.}
\end{equation}
Let \(d_n=\operatorname{lcm}(1,\ldots,2n+2)\).  There are integers
\(A_n,B_n,C_{n,q}\), with \(q\) running over the primes at most \(2n+1\),
such that
\begin{equation}\label{eq:Jn-integer-linear-form}
 d_nJ_n=A_n+B_n\gamma+
 \sum_{\substack{q\le2n+1\\q\ \mathrm{prime}}}C_{n,q}\log q,
 \qquad
 B_n=d_n\binom{2n+1}{n}.
\end{equation}
For every prime \(q\in[n+1,2n+1]\),
\begin{equation}\label{eq:Jn-prime-coefficient-positive}
 C_{n,q}=d_nT_q>0.
\end{equation}
Thus the number of nonzero logarithmic coordinates tends to infinity and is
at least
\(\pi_{\mathrm{pr}}(2n+1)-\pi_{\mathrm{pr}}(n)\sim n/\log n\), where
\(\pi_{\mathrm{pr}}\) is the prime-counting function.

If
\[
 \mathcal H_n=\max\left(
 1,|A_n|,|B_n|,
 \max_{q\le2n+1}|C_{n,q}|
 \right),
\]
then
\begin{align}
 \log\mathcal H_n&=(2+\log4)n+o(n),
 \label{eq:Jn-integer-height}\\
 -\log|d_nJ_n|&=(\log16-2)n+o(n).
 \label{eq:Jn-integer-smallness}
\end{align}
Hence the exponent of this explicit integer form is
\begin{equation}\label{eq:Jn-correct-tau}
 \frac{-\log|d_nJ_n|}{\log\mathcal H_n}
 \longrightarrow
 \frac{\log16-2}{\log4+2}
 =0.2281516725\ldots.
\end{equation}
The exponent is positive, but the ambient dimension grows; this is not a
fixed-dimensional Nesterenko criterion and proves neither irrationality nor
transcendence of \(\gamma\).
\end{theorem}

\begin{proof}
The function
\[
 \phi(x,y)=\log(x(1-x)y(1-y))
\]
has its unique maximum \(-\log16\) at \((1/2,1/2)\), with Hessian
\(-8I_2\).  The remaining factor in \eqref{eq:asymmetric-Jn} has value
\[
 \frac{1-1/2}{(1-1/4)(-\log(1/4))}=\frac1{3\log2}
\]
there.  The two-dimensional Laplace method therefore gives
\[
 J_n\sim16^{-n}\frac{2\pi/n}{8}\frac1{3\log2},
\]
which is \eqref{eq:Jn-correct-asymptotic}.

Equations \eqref{eq:Jn-beta}--\eqref{eq:Jn-alpha-last} show that the
denominators of all \(\alpha_p\) divide
\(\operatorname{lcm}(1,\ldots,n+1)\), while the denominators in
\eqref{eq:Jn-rational-part} divide \(d_n\).  Expanding every
\(\log((p-1)!)\) into prime logarithms proves
\eqref{eq:Jn-integer-linear-form}.  If \(q\in[n+1,2n+1]\) is prime, then
\(2q>2n+1\), so
\[
 v_q((p-1)!)=\mathbf1_{p>q}\qquad(L\le p\le M).
\]
Equation \eqref{eq:Jn-prime-coefficient-positive} follows from
\eqref{eq:Jn-positive-tails}; the prime-number theorem gives the count.

Finally
\[
 \log d_n=2n+o(n),\qquad
 \log\binom{2n+1}{n}=n\log4+o(n).
\]
The coefficient formulas show that every coefficient in
\eqref{eq:Jn-integer-linear-form} is at most
\(d_n4^n\exp(o(n))\), while \(B_n\) has exactly that exponential scale.
This proves \eqref{eq:Jn-integer-height}.  Combining
\eqref{eq:Jn-correct-asymptotic} with \(\log d_n=2n+o(n)\) proves
\eqref{eq:Jn-integer-smallness}, and division gives
\eqref{eq:Jn-correct-tau}.
\end{proof}

\begin{remark}[Comparison with the reciprocal height barrier]
The reciprocal truncations studied earlier saturate their full defining
height so completely that primitive integerization destroys the small
value.  The compact forms \eqref{eq:Jn-integer-linear-form} behave
differently: after integerization they remain exponentially small.  Their
obstruction is instead the rigorously growing set of positive logarithmic
coordinates in \eqref{eq:Jn-prime-coefficient-positive}.  These are two
distinct boundary mechanisms and should not be conflated.
\end{remark}
